% Kapitel 09: Kosmologie, Rotverschiebung und CMB
% Korrigiert: E_H = E0 * xi^(41/4) als Hubble-Energie
% H0 ist SI-Konvention, E_H ist die physikalische Groesse

\chapter{Kosmologie, Rotverschiebung und CMB in der Zeit-Masse-Dualität}

\section{Einführung}

In den vorangegangenen Kapiteln stand die mikroskopische Seite der 
Zeit-Masse-Dualität im Mittelpunkt: Massen, Kopplungen und Quantenphänomene. 
In diesem Kapitel wird gezeigt, wie sich dieselbe Struktur auf großskalige 
Phänomene der Kosmologie auswirkt: Rotverschiebung, kosmische 
Hintergrundstrahlung und effektive Größen wie die Hubble-Skala.

Entscheidend ist: Derselbe Parameter $\xipar = \frac{4}{3} \times 10^{-4}$, 
der auf Teilchenebene die Massen und Kopplungen bestimmt, bestimmt auch die 
kosmologische Rotverschiebung. Es gibt keinen Übergang zu einem anderen 
effektiven Parameter auf großen Skalen -- in einem homogenen, statischen 
Universum gilt überall derselbe $\xipar$.

\section{Rotverschiebung ohne expandierenden Raum}

\subsection{Standard-Interpretation}

Die Standardkosmologie deutet die kosmologische Rotverschiebung hauptsächlich 
als Folge einer expandierenden Raumzeit. Die Wellenlänge eines Photons wird 
mit dem kosmischen Skalenfaktor $a(t)$ mitgedehnt:

\begin{equation}
\frac{\lambda_{\text{obs}}}{\lambda_{\text{emit}}} = \frac{a(t_{\text{obs}})}{a(t_{\text{emit}})} = 1 + z
\end{equation}

\subsection{Zeit-Masse-Dualität Interpretation}

Im Rahmen der Zeit-Masse-Dualität wird die beobachtete Rotverschiebung 
als Folge der fraktalen Tiefenstruktur der Raumzeit verstanden. Ein Photon, 
das durch den fraktalen Raum mit $D_f = 3 - \xipar$ propagiert, verliert 
kontinuierlich Energie an das dynamische Vakuumfeld.

Die T0-Rotverschiebung:

\begin{equation}
z_{\text{T0}} = \int_0^d \xipar(r) \frac{E_\gamma(r)}{E_{\gamma,0}} dr
\end{equation}

Für ein homogenes $\xipar$-Feld vereinfacht sich dies zu:

\begin{equation}
\boxed{z_{\text{T0}} \approx \xipar \cdot d}
\label{eq:z_T0_linear}
\end{equation}

wobei $d$ die Distanz in Megaparsec ist. Diese lineare Beziehung ist die 
T0-Version des Hubble-Gesetzes.

\section{Der Hubble-Parameter und die Hubble-Energie}

\subsection{$H_0$ als SI-Konvention}

$H_0$ wird in km/s/Mpc angegeben. Beim Einsetzen von 
1~Mpc~$= 3{,}26 \times 10^6$~Lj~$= 3{,}08 \times 10^{22}$~m 
kürzt sich Meter heraus und es bleibt eine reine Frequenz:

\begin{equation}
H_0 = 67{,}4\,\text{km/s/Mpc} = 2{,}18 \times 10^{-18}\,\text{s}^{-1}
\end{equation}

Die drei Darstellungen km/s/Mpc, km/s/Lj und s$^{-1}$ sind 
\emph{dasselbe} -- nur verschiedene SI-Einheiten für denselben 
Verhältniswert (Strecke bzw.\ Frequenz).

\subsection{Die Hubble-Energie $E_H$ -- eine neue Größe}

Die Multiplikation einer Frequenz mit $\hbar$ erzeugt eine Energie 
($E = \hbar\nu$, Planck/Einstein). Für $H_0$ definieren wir:

\begin{align}
\text{SI} \to \text{nat.:} \quad & E_H = \hbar \cdot H_0 
\quad\text{(mit $\hbar$ multiplizieren)} \\
\text{nat.} \to \text{SI:} \quad & H_0 = E_H / \hbar
\quad\text{(durch $\hbar$ dividieren)}
\end{align}

$E_H$ ist \emph{nicht mehr $H_0$} (Frequenz), sondern eine Energie. 
In natürlichen Einheiten ($\hbar = 1$) sind $E_H$ und $H_0$ 
numerisch identisch.

\subsection{T0-Vorhersage für $E_H$}

Die T0-Theorie berechnet $E_H$ aus den fundamentalen Parametern 
$\xipar$ und $E_0$:

\begin{equation}
\boxed{E_H = E_0 \cdot \xipar^{41/4} = 7{,}398\,\text{MeV} \cdot \xipar^{10{,}25} = 1{,}41 \times 10^{-33}\,\text{eV}}
\label{eq:E_H_T0}
\end{equation}

Rückrechnung nach SI ($H_0 = E_H / \hbar$):

\begin{equation}
H_0^{\text{T0}} = \frac{E_H}{\hbar} = 2{,}14 \times 10^{-18}\,\text{s}^{-1} = 66{,}2\,\text{km/s/Mpc}
\end{equation}

\subsection{Vergleich mit Messwerten}

\begin{center}
\begin{tabular}{lrr}
\toprule
Messung & $H_0$ [km/s/Mpc] & $E_H$ [eV] \\
\midrule
Planck 2018 & $67{,}4 \pm 0{,}5$ & $1{,}44 \times 10^{-33}$ \\
SH0ES 2022 & $73{,}0 \pm 1{,}0$ & $1{,}56 \times 10^{-33}$ \\
T0 ($E_0 \cdot \xipar^{41/4}$) & $66{,}2$ & $1{,}41 \times 10^{-33}$ \\
\bottomrule
\end{tabular}
\end{center}

Die Abweichung zum Planck-Wert beträgt $-1{,}9\%$. Der Exponent $41/4$ 
bedarf noch einer physikalischen Begründung aus der $\xipar$-Feldtheorie.

\begin{remark}[Modellabhängigkeit des experimentellen $H_0$]
	Der ``experimentelle'' Wert $H_0 \approx 70\,$km/s/Mpc ist keine 
	modellunabhängige Messung. Er wird aus Rohdaten (Rotverschiebungen und 
	Helligkeiten von Standardkerzen) unter der Annahme eines expandierenden 
	Friedmann-Universums extrahiert. Insbesondere:
	\begin{itemize}
	\item Die Luminositätsdistanz $d_L = d \cdot (1+z)$ setzt einen 
	kosmischen Skalenfaktor voraus, den es in T0 nicht gibt.
	\item In T0 ist die physikalische Distanz $d$ direkt messbar, ohne 
	$(1+z)$-Korrektur.
	\item Die modellunabhängige Beobachtung ist die lineare Beziehung 
	$z \propto d$ für kleine $z$. Deren Proportionalitätsfaktor wird 
	in verschiedenen kosmologischen Modellen verschieden interpretiert.
	\end{itemize}
\end{remark}

\section{CMB-Temperatur}

Die CMB-Temperatur:

\begin{equation}
T_{\text{CMB}} = 2.7255\,\text{K}
\end{equation}

wird in der T0-Theorie als thermodynamischer Gleichgewichtszustand der 
$\xipar$-Geometrie interpretiert, nicht als Relikt eines Urknalls. Das 
dynamische Vakuumfeld $\Phi = \rho e^{i\theta}$ hat eine intrinsische 
Phasenentwicklung $\dot{\theta} = m = 1/T$ (aus der Zeit-Masse-Dualität). 
Die CMB-Strahlung ist das thermische Gleichgewichtsspektrum dieses 
universellen Vakuumfelds, dessen Temperatur durch die geometrischen 
Parameter $\xipar$ und $f = 7500$ festgelegt wird.

\section{Statisches Universum im 4D-Torus}

Die T0-Theorie beschreibt ein statisches Universum ohne globale Expansion. 
Die Raumzeit hat die Topologie $\mathbb{R}^3 \times S^1$, wobei $S^1$ die 
Phasenrichtung des Vakuumfelds ist (nicht ``Zeit'' im klassischen Sinn, 
da $T \cdot m = 1$ die Zeit als Kehrwert der Masse definiert).

Das Universum ist ``unendlich'' nicht im Sinne einer unendlichen Ausdehnung, 
sondern weil der Torus in sich geschlossen ist -- es gibt keinen Rand und 
keine Grenze. In dieser Sicht:

\begin{itemize}
\item Rotverschiebung entsteht durch Energieverlust im fraktalen Vakuumfeld, 
nicht durch Expansion
\item Die Hubble-Relation $z = \xipar \cdot d$ ist eine direkte Konsequenz 
der fraktalen Dimension $D_f = 3 - \xipar$
\item Dunkle Energie ist keine separate Substanz, sondern manifestiert sich 
als effektive Eigenschaft des dynamischen Vakuumfelds $\Phi$
\item Die großskalige Homogenität folgt aus der Topologie des Torus, ohne 
Inflation
\end{itemize}

JWST-Beobachtungen entwickelter Galaxien bei $z > 10$, die im 
Standardmodell unerwartet früh erscheinen, sind im T0-Bild natürlich, da 
die Entwicklungszeit unbegrenzt ist.

\section{Zusammenfassung}

Die kosmologischen Vorhersagen der T0-Theorie folgen direkt aus $\xipar$:

\begin{itemize}
\item Rotverschiebung: $z = \xipar \cdot d$ (Energieverlust im Vakuumfeld)
\item Hubble-Energie: $E_H = E_0 \cdot \xipar^{41/4} = 1{,}41 \times 10^{-33}\,$eV 
($-1{,}9\%$ vs Planck)
\item Rückrechnung: $H_0 = E_H / \hbar = 66{,}2\,$km/s/Mpc
\item CMB als Gleichgewichtszustand der Vakuumgeometrie
\item Statisches Universum im 4D-Torus $\mathbb{R}^3 \times S^1$
\end{itemize}

